\documentclass[11pt,a4paper]{article}
\usepackage[utf8]{inputenc}
\usepackage[T1]{fontenc}
\usepackage[brazil]{babel}
\usepackage{amsmath,amssymb,booktabs,geometry}
\geometry{margin=2.5cm}
\title{SPLEEN: formulação matemática do modelo compartimental de perfusão esplâncnica}
\author{Documentação técnica do protótipo}
\date{Versão 0.6.0 -- novembro de 2026}
\begin{document}
\maketitle

\begin{abstract}
Este documento especifica a formulação contínua e os indicadores derivados do protótipo SPLEEN. Trata-se de um modelo compartimental reduzido, exploratório e não validado clinicamente. Os coeficientes são hipóteses computacionais, não parâmetros fisiológicos de referência.
\end{abstract}

\section{Estados, unidades e compartimentos}
Sejam $i\in\{e,int,b,panc\}$ os leitos do estômago, intestino, baço e pâncreas. Os estados são o fluxo arterial regional $Q_i(t)$, em L\,min$^{-1}$, e a pressão portal $P_p(t)$, em mmHg. O vetor de frações basais é
\[
\boldsymbol{\alpha}=(0.18,\,0.62,\,0.13,\,0.07),\qquad
\sum_i\alpha_i=1,
\]
e a condição inicial é $Q_i(0)=1.2\alpha_i$ e $P_p(0)=10$ mmHg.

\section{Entradas e gradiente de perfusão}
Denotam-se por $M$ a pressão arterial média (mmHg), $G$ a pressão de pneumoperitônio (mmHg), $V$ o índice abstrato de suporte vasopressor e $B$ a perda volêmica percentual. A pressão arterial efetiva, a pressão venosa efetiva e o gradiente simplificado são
\begin{align}
M_{\mathrm{eff}} &= \max(20,\,M+0.12V-0.45B),\\
P_{\mathrm{ven}} &= \max(3,\,5+0.20G),\\
\Delta P_{\mathrm{eff}} &= \max(8,\,M_{\mathrm{eff}}-P_{\mathrm{ven}}-0.45G).
\end{align}
Essas expressões são parametrizações fenomenológicas do protótipo, não relações hemodinâmicas empiricamente validadas.

\section{Alvo de fluxo regional}
O fator de retroalimentação portal e o fator de pneumoperitônio intestinal são
\begin{align}
f_p(P_p)&=\max\{0.65,\,1-\gamma_p(P_p-10)\},\\
f_G(G)&=\max\{0.35,\,1-0.025(G-12)\},
\end{align}
onde $\gamma_p$ é o ganho de feedback portal. O fluxo esplâncnico total-alvo é
\[
Q_{\mathrm{tot}}^*=\max\left\{0.12,\,
1.2\frac{\Delta P_{\mathrm{eff}}}{72.2}
(1-0.010B)f_p(P_p)\right\}.
\]
Os pesos regionais não normalizados são
\begin{align}
w_e&=0.18(1+0.0015V), & w_{int}&=0.62f_G(G),\\
w_b&=0.13(1+0.0005V), & w_{panc}&=0.07(1+0.0015V).
\end{align}
Com $W=\sum_iw_i$, os alvos regionais são
\[
Q_i^*=Q_{\mathrm{tot}}^*\frac{w_i}{W}.
\]

\section{Sistema acoplado de EDOs}
A dinâmica dos leitos arteriais é representada por
\begin{equation}
\boxed{\frac{dQ_i}{dt}=\frac{Q_i^*(M,G,V,B,P_p)-Q_i}{\tau_i s_\tau}},
\qquad i\in\{e,int,b,panc\},
\end{equation}
em que $\tau_i$ são constantes de tempo regionais (min) e $s_\tau$ é um fator multiplicativo adimensional. A entrada portal e a saída portal-hepática são
\begin{align}
Q_{\mathrm{portal,in}}&=Q_{int}+Q_b+Q_{panc},\\
Q_{\mathrm{portal,out}}&=\max\left(0,\frac{P_p-P_{hv}}{R_h}\right),
\end{align}
A estrutura atual inclui os fluxos intestinal, esplênico e pancreático na entrada portal. O fluxo do compartimento gástrico integra o fluxo arterial esplâncnico total, mas não é incluído no retorno portal nesta versão reduzida. Essa omissão é uma hipótese estrutural explícita, pendente de revisão anatômica e fisiológica; não deve ser interpretada como ausência de drenagem gástrica para o sistema porta.
\begin{align}
com $P_{hv}=5$ mmHg e $R_h$ em mmHg/(L\,min$^{-1}$). A dinâmica do compartimento portal é
\begin{equation}
\boxed{C_p\frac{dP_p}{dt}=Q_{\mathrm{portal,in}}-Q_{\mathrm{portal,out}}},
\end{equation}
onde $C_p$ é a complacência efetiva (L/mmHg).

\section{Pressão de perfusão abdominal como variável derivada}
A interface calcula a pressão de perfusão abdominal a partir das entradas de pressão arterial média e pressão intra-abdominal:
\[
APP = PAM - PIA.
\]
No estado atual, APP é exibida como indicador derivado, mas o gradiente usado pela função fenomenológica de fluxo-alvo contém termos adicionais para pressão venosa efetiva e pneumoperitônio. Portanto, APP não é ainda a única variável de pressão que dirige as EDOs. A relação entre pressão arterial, pressão intra-abdominal, pressão venosa de saída e resistência vascular requer revisão fisiológica antes de se alterar a formulação.

\section{Indicadores derivados}
O componente arterial hepático simplificado é
\[
Q_{HA}=0.2Q_{\mathrm{tot}}\,
\min\{1.3,\max[0.7,1+0.03(10-P_p)]\},
\qquad Q_{\mathrm{tot}}=\sum_iQ_i.
\]
Para $C_{aO_2}=18.5$ mL/dL, a conversão para mL/L é $10C_{aO_2}=185$ mL/L. A entrega hepática estimada é
\[
D_{O_2,\mathrm{hep}}=(Q_{\mathrm{portal,out}}+Q_{HA})\,185
\quad\mathrm{mL\,O_2\,min^{-1}}.
\]
O fluxo portal de entrada não deve ser confundido com a saída portal-hepática.

\section{Integração numérica}
A referência computacional utiliza Euler explícito:
\begin{align}
Q_i^{n+1}&= \max\left[0,Q_i^n+\Delta t\,
\frac{Q_i^{*,n}-Q_i^n}{\tau_i s_\tau}\right],\\
P_p^{n+1}&=\max\left[P_{hv},P_p^n+\frac{\Delta t}{C_p}
(Q_{\mathrm{portal,in}}^n-Q_{\mathrm{portal,out}}^n)\right].
\end{align}
A implementação deve ser submetida a refinamento de passo temporal, testes de positividade e comparação entre implementações. O truncamento dos estados em seus limites é uma salvaguarda numérica do protótipo e não substitui a análise de estabilidade do sistema contínuo.

\section{Investigações de incerteza e identificabilidade}
A triagem local de identificabilidade usa uma matriz de sensibilidades escaladas
\[
J_{kj}=\frac{p_j}{s_{y_k}}\frac{\partial y_k}{\partial p_j},
\]
aproximada por diferenças finitas, em que $s_{y_k}$ é uma escala de saída. Posto deficiente ou colunas quase colineares indicam efeitos paramétricos pouco distinguíveis naquele cenário e para aquelas saídas. Posto completo não prova identificabilidade estrutural global.

A robustez aos priors é examinada comparando distribuições preditivas obtidas com intervalos uniformes contraídos, de referência e expandidos. Essas distribuições são cenários de estresse de hipótese, não priors empíricos.

O experimento sintético de recuperação gera observações a partir de parâmetros conhecidos e adiciona ruído Gaussiano. A aproximação posterior por amostragem por importância permite calcular ranks e erros de recuperação. O diagnóstico é preliminar: tamanho efetivo baixo da amostragem por importância indica degeneração e impede interpretação de calibração.

\section{Limitações e uso}
O modelo não representa microcirculação, autorregulação validada, controle neural/hormonal, retorno venoso completo, farmacocinética, heterogeneidade espacial ou anatomia individual. Os coeficientes são ilustrativos. Não há calibração ou validação independente. O SPLEEN destina-se exclusivamente a pesquisa metodológica e educação; não deve orientar diagnóstico, prescrição ou decisão intraoperatória.

\end{document}
